\subsection{平衡的稳定性}

受力平衡方程只能确定可能的平衡位置；要判断物体受到微小扰动后能否回到原位置，还要
研究平衡的稳定性。如下图所示，同样的两根竖直杆，第一根由顶部悬挂，第二根以下端支承
在水平桌面上。理想情况下，两者都能满足静力平衡条件，但前者是稳定平衡，后者是
不稳定平衡。

\sourcefigure[0.69\linewidth]{fig-1-4.png}
\origcaption{图 1.4：稳定和不稳定平衡}

如果系统稍微偏离平衡位置后出现使偏离量减小的回复力或回复力矩，该平衡就是稳定的。
例如悬挂杆偏离竖直方向后，重力矩使偏角减小。反之，如果扰动后的力或力矩使偏离量
继续增大，该平衡就是不稳定的；竖直支在桌面上的杆正是这种情形。对保守系统，也可以
用势能判断：稳定平衡对应势能的严格局部极小值，不稳定平衡通常对应局部极大值。

不稳定平衡仍然是严格意义上的平衡，并非“不可能出现”；只是它不能抵抗微小扰动。
这类平衡在倒立摆、结构屈曲和临界状态分析中都有重要意义。工程应用中必须把“满足平衡
方程”和“具有稳定性”分开检验；稳定性既可以定性判断，也可以通过力矩展开或势能展开
作定量分析。

\begin{origexamplebox}{例 1.28}
一根质量为 $m$ 的均匀杆，长为 $L$，下端可绕固定水平轴转动。有两根水平弹簧，
劲度系数相同，拴在杆的上端，使其处于竖直位置。问：弹簧的劲度系数 $k$ 为何值，才能使
杆处于稳定平衡？两弹簧按图对称布置，外端固定在竖直杆顶点两侧等高、等距的位置；
每根弹簧的原长记为 $\ell_0>0$，初始长度记为 $d_0>0$。按理想轻弹簧模型讨论小幅平面偏转。
\sourcefigure[0.49\linewidth]{ex-1-28.png}
\end{origexamplebox}

\begin{xesolution}
设杆从竖直方向偏转一个小角 $\varphi$。杆上端的水平位移为
\[
  x=L\sin\varphi=L\varphi+O(\varphi^3).
\]
\Needspace{66mm}
偏转后的杆、重力与两弹簧形变如解图所示：
\soldiagOneTwentyEight
在竖直位置两弹簧关于杆对称。由解图可见，小角度下左、右弹簧的长度改变量一阶项分别为
$+x$ 和 $-x$。因而两弹簧的总弹性势能相对竖直位置的变化为
\[
  \Delta U_{\rm s}=kL^2\varphi^2+O(\varphi^4).
\]
杆的重力势能变化为
\[
  \Delta U_{\rm g}=mg\frac L2(\cos\varphi-1)
  =-\frac{mgL}{4}\varphi^2+O(\varphi^4).
\]
因此总势能在 $\varphi=0$ 附近满足
\[
  U=U(0)+\left(kL^2-\frac{mgL}{4}\right)\varphi^2+O(\varphi^4).
\]
二次项系数为正时，势能在竖直位置为严格局部极小值，有线性回复力矩。因此
\[
  \boxed{k>\frac{mg}{4L}}.
\]
若 $k<mg/(4L)$，二次项为负，竖直位置不稳定。等号处二次项消失，必须继续检查，不能只据二次项下结论。
杆顶下降量为 $v=L(1-\cos\varphi)$，根据两固定点到杆顶的距离，弹簧长度为
\[
 \ell_\pm=\sqrt{(d_0\pm x)^2+v^2}.
\]
将两根弹簧的势能分别写成 $k(\ell_\pm-\ell_0)^2/2$，减去初态势能，有
\[
 \Delta U_{\rm s}
 =k\left[x^2+v^2-\ell_0(\ell_++\ell_--2d_0)\right].
\]
这里用了 $\ell_+^2+\ell_-^2=2d_0^2+2x^2+2v^2$。再利用
$x^2+v^2=2L^2(1-\cos\varphi)$，加上重力势能，得精确关系
\[
 \Delta U=\left(2kL^2-\frac{mgL}{2}\right)(1-\cos\varphi)
 -k\ell_0(\ell_++\ell_--2d_0).
\]
当 $k=mg/(4L)$ 时，第一项恰好消失。只要 $\varphi\ne0$，杆顶就低于两固定点的连线；
三角形两边长度之和大于第三边，故 $\ell_++\ell_->2d_0$。由于 $\ell_0>0$，此时
$\Delta U<0$，所以临界点仍不稳定，稳定范围确为上面的严格不等式。
若另取原长恰为零的理想极限，则临界时 $\Delta U\equiv0$，为随遇平衡，而不是有回复力的稳定平衡。
\end{xesolution}

\begin{origexamplebox}{例 1.29}
两根长度均为 $L$、质量忽略不计的刚性轻杆固定在竖直平面内，顶端由铰链连接
并固定有一个质量为 $M$ 的质点。两杆下端置于同一光滑水平支承面上，可以水平滑动；
下端之间连接一根劲度系数为 $k$、原长为 0 的轻弹簧。求平衡时杆与垂线的夹角 $\theta$。
\sourcefigure[0.39\linewidth]{ex-1-29.png}
\end{origexamplebox}

\begin{xesolution}
题图中的光滑水平支承面只提供竖直支持力，$A$、$B$ 可以沿水平方向滑动。
解图分别画顶点 $M$ 和下端 $A$ 的受力，其中 $N_A$ 是支承面对 $A$ 的竖直支持力；
$B$ 端按对称性处理，$N_B=N_A$：
\soldiagOneTwentyNine

弹簧原长为零，平衡时长度为 $2L\sin\theta$，所以弹簧张力
\[
  T=2kL\sin\theta.
\]
设每根杆中的压力为 $P$。顶端质点的竖直平衡给出
\[
  2P\cos\theta=Mg.
\]
在下端 $A$ 或 $B$ 处列水平方向平衡方程：
\[
  P\sin\theta=T=2kL\sin\theta.
\]
因此完整的平衡条件可以写成
\[
  \sin\theta\bigl(4kL\cos\theta-Mg\bigr)=0.
\]
它包含两个分支。第一支是 $\theta=0$，此时 $A$、$B$ 重合、两杆重叠，是图示构型的退化
极限。第二支为非退化构型 $\sin\theta\ne0$，此时 $P=2kL$，并有
\[
  \boxed{\cos\theta=\frac{Mg}{4kL}},
  \qquad
  \boxed{\theta=\arccos\frac{Mg}{4kL}}.
\]
非退化分支存在的条件为
\[
  \boxed{k>\frac{Mg}{4L}}.
\]

由于本题位于“平衡的稳定性”一节，还应继续检查稳定性。在上述补充模型下，系统势能为
\[
  U(\theta)=MgL\cos\theta+2kL^2\sin^2\theta.
\]
在非退化平衡角处
\[
  U''(\theta)=-4kL^2\sin^2\theta<0,
\]
所以方框中的非零夹角分支是\textbf{不稳定平衡}。退化的 $\theta=0$ 分支在
$k>Mg/(4L)$ 时稳定；等号处二次项消失而四次项为负，仍不稳定。
\end{xesolution}

\begin{origexamplebox}{例 1.30}
一根长为 $L$、质量为 $m$ 的均匀杆一端靠于竖直墙面，另一端靠在曲面上，墙面与曲面均光滑。
取水平向右为 $x$ 轴、竖直向上为 $y$ 轴，曲线过原点。要杆在图示构型的每个可实现位置都平衡，
求曲面的截面满足的条件，并说明适用的位置范围。
\sourcefigure[0.42\linewidth]{ex-1-30.png}
\end{origexamplebox}

\begin{xesolution}
设杆与曲面的接触点为 $P(x,y)$，杆在墙上的端点为 $A(0,Y)$。由杆长为 $L$，
\[
  Y-y=\sqrt{L^2-x^2}.
\]
三个力的作用线及其交点 $C$ 如解图所示：
\soldiagOneThirty
假设墙面和曲面均光滑。杆只受墙的水平支持力、曲面的法向支持力和重力三个力。三力
平衡时，其作用线必须交于一点。墙面支持力的作用线是过 $A$ 的水平线，重力作用线是过
杆中点的竖直线，两者交点为
\[
  C\left(\frac x2,Y\right).
\]
所以曲面在 $P$ 点的法线必须通过 $C$。直线 $PC$ 的斜率为
\[
  \frac{Y-y}{x/2-x}=-\frac{2\sqrt{L^2-x^2}}{x}.
\]
另一方面，曲线 $y(x)$ 的法线斜率为 $-1/y'$，于是
\[
  y'=\frac{x}{2\sqrt{L^2-x^2}}.
\]
利用 $\mathrm d\sqrt{L^2-x^2}/\mathrm dx=-x/\sqrt{L^2-x^2}$，有
\[
 y(x)=-\frac12\sqrt{L^2-x^2}+C_0.
\]
曲线过原点，代入 $x=0,y=0$ 得 $C_0=L/2$，于是
\[
  \boxed{y(x)=\frac{L-\sqrt{L^2-x^2}}{2}}.
\]
等价地，曲线满足
\[
  \boxed{x^2+4\left(y-\frac L2\right)^2=L^2},
\]
即它是半长轴为 $L$、半短轴为 $L/2$ 的椭圆下半支中位于墙右侧的一段。
杆的质心高度为
\[
 y_G=\frac{Y+y}{2}=y+\frac12\sqrt{L^2-x^2}=\frac L2,
\]
不随位置改变，因此在可保持接触的位置上为随遇平衡。

还需检验接触力和端点。曲面法线的水平、竖直分量之比为 $-y':1$；由竖直平衡和水平平衡，
\[
 N_{P,y}=mg,\qquad
 N_A=mg\,y'=\frac{mgx}{2\sqrt{L^2-x^2}},\qquad
 N_P=mg\sqrt{1+y'^2}.
\]
在 $0<x<L$ 内两处支持力均不要求拉力。$x=0$ 是竖直杆极限，此时 $N_A=0$、$N_P=mg$，
可单独由受力平衡验证，不能在前面含 $1/x$ 的斜率式中直接代入零。
当 $x\to L^-$ 时，$N_A$ 和 $N_P$ 都发散；在 $x=L$ 处杆水平，两处接触法线均水平，
不能平衡重力。因此范围是
\[
 \boxed{0\le x<L},
\]
不包括最右端点。实际材料强度有限时，可用范围还会比这一理想范围更小。
\end{xesolution}

\begin{origexamplebox}{例 1.31【思考题】}
\begin{minipage}[t]{0.62\linewidth}
如图，一个固定圆形“碗”的半径为 $R$，一个小球半径为 $r$；接触面足够粗糙，小球
作纯滚动而不滑动，初始时小球重心位于球心正上方。小球不是匀质
的，设 $R>r$，且重心在球内、$0<h=OG<r$。请问 $h$ 为多大时形成稳定平衡。如果半径 $R$ 不断增大，会怎样？
增大到无穷大会怎样？再进一步增大会怎样？
\end{minipage}\hfill
\begin{minipage}[t]{0.34\linewidth}
\vspace{0pt}\centering
\resizebox{\linewidth}{!}{\faithfulExOneThirtyOne}
\end{minipage}
\end{origexamplebox}

\begin{xesolution}
设 $OG=h$，并取 $h>0$ 表示重心位于球心上方。假设小球在碗内滚动而不滑动，记
$a=R-r$。小球偏离最低位置后的两个转角和重心位置如解图所示：
\soldiagOneThirtyOne
球心从最低点沿碗转过小角 $\varphi$ 时，平动弧长为 $a\varphi$。无滑动条件给出小球
相对地面的转角
\[
  \psi=-\frac{a}{r}\varphi.
\]

略去与 $\varphi$ 无关的常数后，重心的高度可以写成
\[
  y_G=-a\cos\varphi+h\cos\frac{a\varphi}{r}.
\]
在 $\varphi=0$ 附近展开：
\[
  y_G=y_G(0)+\frac12\left(a-\frac{ha^2}{r^2}\right)\varphi^2
  +O(\varphi^4).
\]
二次项系数为正时，最低位置是势能的严格局部极小点，具有线性回复力。因此
\[
  \boxed{h<\frac{r^2}{R-r}}.
\]
当 $h>r^2/(R-r)$ 时二次项为负，最低位置不稳定。等号不能只凭二次项判断。
展开两个余弦至四阶：
\[
 y_G-y_G(0)=\frac12\left(a-\frac{ha^2}{r^2}\right)\varphi^2
 +\frac1{24}\left(-a+\frac{ha^4}{r^4}\right)\varphi^4+O(\varphi^6).
\]
在 $h=r^2/a$ 处，四次项系数化为
\[
 \frac1{24}\left(-a+\frac{a^3}{r^2}\right)
 =\frac{a(a^2-r^2)}{24r^2}.
\]
$R>2r$ 时 $a>r$，该系数为正，临界点仍是严格势能极小点，只是回复力矩从三次项开始，
称为四阶弱稳定。$R=2r$ 时 $a=r,h=r$，两项余弦完全抵消，为随遇平衡；
$r<R<2r$ 时形式上的临界四次项为负，但所需 $h=r^2/a>r$ 已不可能位于球内。

结合题设 $0<h<r$，物理上可实现的稳定范围为
\[
 \boxed{\begin{cases}
  0<h<r, & r<R\le2r,\\
  0<h\le r^2/(R-r), & R>2r.
 \end{cases}}
\]
第二行的等号包含弱稳定；若专指有非零线性回复力的稳定，则把等号去掉。
例如 $R=3r,h=r/2$ 恰是可实现的弱稳定临界状态。

当 $R$ 增大时，$h_{\rm c}=r^2/(R-r)$ 单调减小；$R\to\infty$ 时碗面趋于水平面，
$h_{\rm c}\to0$，任何固定的 $h>0$ 最终都会变得不稳定。平面上球心高度不变，
重心向上时滚动会使其降低，所以这个极限也可直接理解。

“再进一步增大”不是把无穷大当作普通数继续相加，而是让表面的曲率从凹面经零曲率变成凸面。
下面另取凸面半径 $S>0$，仍令 $h>0$ 表示初态重心在球心上方，$h<0$ 表示在下方。
\Needspace{65mm}
凸面圆心、接触点和小球球心的关系如解图所示：
\soldiagConvexRolling
记球心轨道半径为 $a_+=S+r$，取球心向右偏转的角度为 $\varphi$。
球心速度沿凸面切线向右下方；接触点无滑动要求小球顺时针转动，故
\[
 \psi=-\frac{a_+}{r}\varphi.
\]
以凸面圆心为高度原点，球心高度为 $a_+\cos\varphi$，重心相对球心的竖直分量为
$h\cos(a_+\varphi/r)$，因而
\[
 y_G=a_+\cos\varphi+h\cos\frac{a_+\varphi}{r},
\]
\[
 y_G-y_G(0)=-\frac12\left(a_++\frac{ha_+^2}{r^2}\right)\varphi^2+O(\varphi^4).
\]
二次项为正给出
\[
 \boxed{h<-\frac{r^2}{S+r}}.
\]
即重心必须位于球心下方，且下移足够多。在等号处，四次项系数为
$a_+(1-a_+^2/r^2)/24<0$，所以临界点不稳定，等号不取。
该条件还需与重心在球内的限制 $-r<h<r$ 同时满足。
\end{xesolution}
